% Methods section

\subsection{DFT Calculations}

All density-functional theory (DFT) calculations were performed using VASP~5.4.4~\cite{vasp_a,vasp_b} with the projector augmented wave (PAW) method~\cite{paw_blochl,paw_kj} and the Perdew--Burke--Ernzerhof (PBE) exchange-correlation functional~\cite{pbe}. PAW potentials from the VASP recommended-PBE set were used: \texttt{Fe\_pv} and \texttt{Mo\_sv}. The complete set of INCAR parameters is reproduced from the high-throughput workflow file (\texttt{vasp\_tight.yaml}, deposited on Zenodo~\cite{zenodo}) in Table~\ref{tab:dft_settings}.

\begin{table}[htbp]
  \centering
  \caption{DFT settings used for all 332 calculations (training and validation sets). Parameters are taken verbatim from the workflow file \texttt{vasp\_tight.yaml}.}
  \label{tab:dft_settings}
  \begin{tabular}{ll}
    \toprule
    Parameter & Value \\
    \midrule
    ENCUT  & 500~eV \\
    PREC   & Accurate \\
    EDIFF  & $10^{-6}$~eV \\
    ISMEAR & 0 (Gaussian smearing) \\
    SIGMA  & 0.1--0.2~eV$^{a}$ \\
    ISPIN  & 1 (non-spin-polarised) \\
    NELM / NELMIN & 500 / 20 \\
    LREAL  & Auto \\
    $k$-spacing & 0.125~\AA$^{-1}$ (Monkhorst--Pack) \\
    \bottomrule
  \end{tabular}
  \begin{flushleft}
    \small $^{a}$ The workflow YAML specifies \texttt{sigma: 0.2}; the run-set directory tag encodes 0.1. Both values are reported as the inconsistency cannot be resolved without the original OUTCAR archives. The entropy correction $T \cdot S$ is $<$1~\meVat for $\sigma \leq 0.2$~eV at 0~K extrapolation.
  \end{flushleft}
\end{table}

All structures, including the large complex TCP cells (R, M, P, $\delta$), were sampled with $\Gamma$-centred Monkhorst--Pack meshes at fixed reciprocal-space density 0.125~\AA$^{-1}$ as specified in the workflow. Table~S1 provides a sensitivity check confirming that $\Gamma$-only and $2\times2\times2$ formation energies agree to within 0.3~\meVat (the EOS-fit residual level); the fixed-density MP protocol was used in production.

For each structure, a series of constant-volume calculations spanning a range of volumes around equilibrium was performed, followed by $E$--$V$ curve fitting with a third-order Birch--Murnaghan EOS~\cite{birch_murnaghan}. This eliminates Pulay stress errors. The dataset curation (300 raw calculations $\to$ 291 $\to$ 262 TCP structures) is documented in Supplementary~\S S1; per-structure BM-fit metrics are tabulated in the deposited \texttt{FullyCuratedParsedBriefSummary.json}~\cite{zenodo}.

The target property for machine learning was the formation energy referenced to non-magnetic (NM) hcp Fe and hcp Mo, denoted $E_f^{nmhcp}$. All calculations enforced non-spin-polarised occupancies (\texttt{ISPIN=1}), consistent with the paramagnetic behaviour of Fe--Mo TCP phases at the compositions and temperatures studied. A spin-polarised cross-check on Fe-rich $\chi$ and $\mu$ configurations was not performed in this study and is identified as a known limitation; for reference, the local moments on Wyckoff sites in those phases have been reported in the literature but are small enough that the relative ordering of TCP polymorph energies is not expected to change~\cite{smith_tcp}. The magnetic energy of bcc Fe relative to NM hcp Fe is approximately 0.5~eV/atom, which shifts absolute formation energies but does not affect the relative energy differences between TCP polymorphs. Training (262) and validation (70) structures were generated by the same \texttt{ams\_highthroughput} workflow invocation using a single \texttt{vasp\_tight.yaml}, guaranteeing identical DFT settings across both sets.

\subsection{Descriptor Computation}

\subsubsection*{BOP Moments}

Bond-order potential (BOP) moments were computed using the BOPfox code~\cite{bopfox}. The production descriptor used the \textbf{\texttt{0.7projections\_os}} parameterisation: an orthogonal $d$-only valence model with on-site levels rescaled by 0.5 (Fe: $-3.61$~eV, Mo: $-1.68$~eV) and bond integrals ($dd\sigma$, $dd\pi$, $dd\delta$) from DFT$\to$TB projections (sum7exp form) rescaled by 0.7. The BOPfox keys used were: \texttt{moments~=~20}, \texttt{terminator~=~constantabn}, \texttt{bandwidth~=~findeminemax}, \texttt{bopkernel~=~jackson}, \texttt{nexpmoments~=~200}. Twenty moments per atom were retained. The full parameter file (\texttt{FeMo\_0.7projections\_os.bx}) is deposited on Zenodo~\cite{zenodo}. A canonical $d$-band TB variant (\texttt{Fe-Mo.bx}, power-law bond integrals, 9 moments) was used only as an ablation; all headline results are from the projection model. After CNAV aggregation, BOP feature vectors had approximately 405--417 dimensions, with the exact count depending on the number of CN groups in each phase.

\textbf{Note on BOPfox availability.} BOPfox is not openly distributed; it is available from the authors of that code upon reasonable request. To ensure that the downstream regression is reproducible without BOPfox, we deposit the pre-computed BOP descriptor matrix for all 262$+$70 structures (\texttt{parallel\_Fe-Mo\_relaxed\_0.7projections\_0.5os\_table\_WUBIND\_20.pkl}) and the corresponding CNAV CSVs on Zenodo~\cite{zenodo}.

\subsubsection*{ACE Descriptors}

Atomic cluster expansion (ACE) descriptors~\cite{ace} were computed using the \texttt{python-ace} package (installed from source with \texttt{python-ace-cmake-compat.patch}; requires Cython~$<$~3). We used a cutoff radius of \SI{6.0}{\angstrom}, correlation orders 2--4, and generic radial basis functions spaced at \SI{0.5}{\angstrom} intervals. The full ACE basis yielded up to 1803 features per structure. A ``canonical'' variant using only the dominant radial basis functions reduced this to 417 features while retaining most predictive accuracy. Note that ACE was used here as a fixed structural fingerprint, not as a fitted interatomic potential; a comparison with a potential fitted via \texttt{pacemaker} on the same dataset would require the original OUTCAR forces and stresses, which are no longer available (see \S\ref{sec:limitations}).

\subsubsection*{SOAP Descriptors}

Smooth overlap of atomic positions (SOAP) descriptors~\cite{soap} were computed using the DScribe library (v2.1.2)~\cite{dscribe}. The SOAP power spectrum was calculated with: $r_\text{cut} = \SI{6.0}{\angstrom}$, $n_\text{max} = 8$ radial basis functions, and $l_\text{max} = 6$ spherical harmonics. Element-specific SOAP vectors (Fe--Fe, Fe--Mo, Mo--Mo, Mo--Fe) were computed for each central atom and concatenated. The resulting per-atom SOAP vectors had dimension~287 after CNAV aggregation.

\subsubsection*{CNAV Aggregation}

For each structure, the per-atom descriptor vectors were averaged within groups sharing the same coordination number (CN). The CN for each atom was determined by a hard cutoff at \SI{3.5}{\angstrom} on the first shell; this distance falls well within the first-neighbour gap of Fe--Mo TCP phases ($\sim$2.5--2.9~\AA{} first shell, next shell $>$4.0~\AA), making CN-bin assignment robust to small cutoff variations. The CN-resolved averages were concatenated into a single structure-level fingerprint, with the number of CN groups varying between 3 and 6 depending on the phase.

Crucially, the geometric CN groups obtained with this cutoff coincide with the Frank--Kasper polyhedral classification (CN12, CN14, CN15, CN16)~\cite{frank_kasper_1,frank_kasper_2} for all inequivalent Wyckoff sites of the R, M, P, and $\delta$ phases, so CNAV groups inherit the Frank--Kasper site hierarchy rather than averaging across it. Supplementary Table~S2 provides the full Wyckoff $\leftrightarrow$ Frank--Kasper $\leftrightarrow$ CNAV-group mapping for each complex TCP phase, derived from the prototype structure files used in the pipeline (\texttt{POSCAR\_\{R,M,P,delta\}\_proto.vasp}).

CNAV preserves crystallographic information by treating each CN environment separately, while reducing noise from per-atom fluctuations. For BOP and SOAP descriptors, CNAV reduces test-set RMSE by 20--30\% relative to per-structure averaging; for ACE descriptors the improvement is within numerical noise on the present single-seed split, suggesting that ACE already encodes CN-resolved environments more explicitly (Supplementary Table~S3). The ensemble gives 5--15\% lower RMSE than the best individual model per descriptor family.

\subsection{Model Training and Evaluation}

All machine learning models were implemented using scikit-learn~v1.3.0~\cite{sklearn} with Python~3.10.20, ASE~3.22.1, and NumPy~1.26.4; the full conda environment lockfile is deposited on Zenodo~\cite{zenodo}. The training/test split (209/53) was stratified by phase to ensure that each simple TCP phase was represented in both sets. Features were standardised to zero mean and unit variance using training-set statistics. Three seeds were used: DFT workflow seed 42 (\texttt{setup.yaml}); feature-selection seed 20091116 (\texttt{Scripts/FeatureSelection.py}); recursivity-search seed 23192 (\texttt{Scripts/TestRecursivity.py}).

\subsubsection*{Model Selection}

Three regression algorithms were considered:

\begin{enumerate}
  \item \textbf{Kernel Ridge Regression (KRR):} \emph{Polynomial} kernel (degrees 2--5, \texttt{coef0} 0--5), regularisation $\alpha \in \{1, 0.1, 0.01\}$, optimised via \texttt{GridSearchCV} with 5-fold stratified CV on the 209-structure training partition.
  \item \textbf{Random Forest (RF):} 200 trees; \texttt{max\_depth} $\in \{5, 10\}$, \texttt{min\_samples\_leaf} $\in \{1, 10\}$, \texttt{max\_leaf\_nodes} $\in \{\text{None}, 20, 30\}$; grid search with 5-fold CV.
  \item \textbf{Multi-Layer Perceptron (MLP):} sklearn \texttt{MLPRegressor}; hidden layer sizes $\in \{(20,4), (40,4)\}$; \texttt{solver~=~lbfgs}; \texttt{activation~=~logistic}; $\alpha \in \{0.03, 0.04, 0.05, 0.06, 0.1\}$; \texttt{random\_state~=~20091116}; grid search with 5-fold CV.
\end{enumerate}

\subsubsection*{Ensemble Construction}

For each feature set, a VotingRegressor ensemble was constructed by averaging the predictions of the three optimally tuned models (KRR, RF, MLP). This approach consistently outperformed individual models, with the ensemble RMSE being 5--15\% lower than the best individual model.

\subsubsection*{Feature Selection}

Feature selection used \texttt{RFECV} (sklearn recursive feature elimination with cross-validation) with a 5-fold \texttt{StratifiedKFold} over the \emph{full} 262-structure dataset. Consequently, the 53-structure held-out test set was visible to the feature selector; Table~\ref{tab:test_performance} test-set RMSEs are therefore optimistic estimates, and the fully independent 70-structure complex-TCP validation set (Table~\ref{tab:validation_rmse}) is the primary measure of generalisation. Feature selection was run 10--13 times per descriptor family (embarrassingly parallel, 16 cores per job); the distribution of selected feature counts across runs is reported in Supplementary Table~S4 (median and IQR). Starting from the full feature set, features were eliminated based on 5-fold CV RMSE until no further improvement exceeded 0.5~\meVat, reducing dimensionality by factors of 2--10 (from $\sim$400--1800 to $\sim$50--300). Hyperparameter grids are deposited in \texttt{Fe-Mo/results/\{Kernel Ridge|MLP|Random Forest\}\_options.json}~\cite{zenodo}.

\subsubsection*{Evaluation Metrics}

Model performance was evaluated using:
\begin{align}
  \text{RMSE} &= \sqrt{\frac{1}{N} \sum_{i=1}^{N} (y_i - \hat{y}_i)^2}, \\
  \text{MAE} &= \frac{1}{N} \sum_{i=1}^{N} |y_i - \hat{y}_i|.
\end{align}
Both metrics were computed on the test set (53 structures) and the independent validation set (70 structures).

\subsubsection*{Statistical Reporting}

All test-set and validation-set RMSE values reported in Tables~\ref{tab:test_performance}--\ref{tab:validation_rmse} and in the text are \textbf{point estimates} from a single fitted VotingRegressor ensemble on the stratified 80/20 train/test split with seed~=~42. Confidence intervals are not provided because the raw descriptor matrices required for re-bootstrapping are no longer available. The only confidence interval explicitly computed in the pipeline is the 95\% percentile band on the parity-plot regression line in Fig.~\ref{fig:validation}, generated by a $B=1000$ non-paired bootstrap of $(y_\text{true}, y_\text{pred})$ pairs (\texttt{11\_ValidatePredictions.ipynb}, cell~15). No formal pairwise hypothesis test (paired $t$, Wilcoxon, or paired bootstrap of residuals) was performed; the BOP~$<$~ACE~$<$~SOAP ordering reported in the text is a descriptive observation consistent across all four complex TCP phases and across all three constituent models of the VotingRegressor. Supplementary Table~S4 lists per-run selected feature counts (10--13 runs per descriptor) extracted from preserved notebook outputs.

\subsection{Thermodynamic Modelling}

Finite-temperature thermodynamic properties were computed within the compound energy formalism (CEF) using the Bragg--Williams (BW) approximation, as implemented in the \texttt{pycef} code. The Gibbs free energy of each phase is given by:

\begin{equation}
  G = \sum_i y_i G_i^\circ + RT \sum_s \sum_i y_i^{(s)} \ln y_i^{(s)} + G^E
  \label{eq:gibbs}
\end{equation}

where $y_i^{(s)}$ is the site fraction of species $i$ on sublattice $s$, $G_i^\circ$ is the ML-predicted formation energy of the end-member compound, and $G^E$ is the excess Gibbs free energy. The ideal BW approximation ($G^E = 0$) was used as a first approximation, neglecting short-range order (SRO). To test the sensitivity to this assumption, we performed calculations with small excess terms $L = \pm 5$ and $\pm 10$~kJ/mol. The R-phase site occupancies changed by less than 0.03 in site fraction for all but the mixed 6c~site ($\Delta y_{\text{Mo}} \sim 0.08$), justifying the ideal BW approximation at the accuracy level of this study (Supplementary Figure~S2). The calculation was performed at $T = 1700$~K, at which temperature the Gibbs free energy curves of all studied TCP phases can be compared on a common axis. We note, however, that this temperature lies \emph{above} the upper stability bound of the R phase in the assessed Fe--Mo diagram: Houserová \emph{et al.}~\cite{houserova_2002} place $T_\text{max}^\text{R} \approx 1488$~K and Andersson~\cite{andersson_1988} places it at $\approx 1500$~K. The R-phase sublattice-occupancy comparison in \S\ref{sec:thermo} should therefore be interpreted as a \emph{trend-level} test of the ML-predicted end-member ordering, not as a strict thermodynamic equilibrium comparison at the experimental annealing temperature; recomputing at $T = 1473$~K (the annealing temperature of the Joubert 1993 refinement~\cite{joubert_fe_mo}) is identified as follow-up work. Sublattice occupancies were obtained by minimising Eq.~\eqref{eq:gibbs} with respect to site fractions using a sequential least-squares programming (SLSQP) optimiser~\cite{slsqp} from 10 random starting points; the same minimum was found in all cases (empirical confirmation, not a proof of global optimality).

\subsection{Data and Code Availability}

Pre-computed descriptors (BOP descriptor matrix, CNAV CSVs, ACE and SOAP feature tables), fitted model pickles, hyperparameter JSON files, prediction CSVs, the workflow inputs (\texttt{vasp\_tight.yaml}, \texttt{setup.yaml}), the BOPfox parameter file (\texttt{FeMo\_0.7projections\_os.bx}), and the conda environment lockfile are deposited in the Zenodo repository (DOI: \texttt{10.5281/zenodo.19427673})~\cite{zenodo}. The full pipeline code and Jupyter notebooks are available on GitHub at \url{https://github.com/AIIMProject/MLFeMoTCPs}. The raw VASP OUTCAR archives are no longer available to the authors; the deposited BOP descriptor matrix and CNAV CSVs allow the downstream ML regression to be reproduced bit-identically without BOPfox. The BOPfox code itself is available from the developers upon reasonable request; the deposited \texttt{.bx} parameter file makes the descriptor pipeline fully specified for any user with BOPfox access.
